% ============================================================
%  Chapitre 3 — Logique Modale
% ============================================================
\chapter{Logique Modale}

\section{Fondements théoriques}

\subsection{Syntaxe}

La \textbf{logique modale} étend la logique propositionnelle par deux opérateurs
unaires :

\begin{itemize}
  \item $\nec\varphi$ (\emph{nécessairement} $\varphi$) : $\varphi$ est vraie
        dans \emph{tous} les mondes accessibles.
  \item $\pos\varphi$ (\emph{possiblement} $\varphi$) : $\varphi$ est vraie
        dans \emph{au moins un} monde accessible.
\end{itemize}

\noindent La grammaire des formules modales est :
\[
  \varphi \;::=\; p \;\mid\; \lnot\varphi \;\mid\; \varphi \land \varphi
           \;\mid\; \varphi \lor \varphi \;\mid\; \varphi \to \varphi
           \;\mid\; \nec\varphi \;\mid\; \pos\varphi
\]
où $p$ est un atome propositionnel.

\subsection{Sémantique — Modèles de Kripke}

Un \textbf{modèle de Kripke} est un triplet $\mathcal{M} = \langle W, R, V \rangle$ :

\begin{itemize}
  \item $W$ : ensemble non vide de \emph{mondes possibles} ;
  \item $R \subseteq W \times W$ : \emph{relation d'accessibilité} entre mondes ;
  \item $V : \mathrm{Prop} \to 2^W$ : fonction de \emph{valuation} associant à
        chaque atome l'ensemble des mondes où il est vrai.
\end{itemize}

\noindent La relation de satisfaction $\mathcal{M}, w \models \varphi$ est définie
inductivement ; les cas modaux sont :

\begin{align*}
  \mathcal{M}, w &\models \nec\varphi
    &&\iff \forall v \in W,\; (w, v) \in R \Rightarrow \mathcal{M}, v \models \varphi \\
  \mathcal{M}, w &\models \pos\varphi
    &&\iff \exists v \in W,\; (w, v) \in R \land \mathcal{M}, v \models \varphi
\end{align*}

\subsection{Propriétés de la relation d'accessibilité}

Les propriétés de $R$ définissent les systèmes modaux (Tableau~\ref{tab:modal-systems}).

\begin{table}[H]
\centering
\caption{Propriétés de $R$ et systèmes modaux correspondants}
\label{tab:modal-systems}
\begin{tabular}{llll}
\toprule
\textbf{Propriété} & \textbf{Définition} & \textbf{Axiome} & \textbf{Système} \\
\midrule
Réflexive   & $\forall w,\; (w,w) \in R$               & $\nec p \to p$        & \textbf{T}  \\
Transitive  & $(w,v),(v,u) \in R \Rightarrow (w,u) \in R$ & $\nec p \to \nec\nec p$ & \textbf{S4} \\
Symétrique  & $(w,v) \in R \Rightarrow (v,w) \in R$   & $p \to \nec\pos p$    & \textbf{B}  \\
Euclidienne & $(w,v),(w,u) \in R \Rightarrow (v,u) \in R$ & $\pos p \to \nec\pos p$ & \textbf{S5} \\
Sérielle    & $\forall w,\; \exists v,\; (w,v) \in R$  & $\nec p \to \pos p$   & \textbf{D}  \\
\bottomrule
\end{tabular}
\end{table}

\section{Application — Industrie musicale}

\subsection{Modèle de Kripke}

\[
  W = \{w_1=\text{label},\; w_2=\text{fan},\; w_3=\text{artiste},\; w_4=\text{critique}\}
\]
\[
  R = \{(w_1,w_2),\,(w_1,w_3),\,(w_3,w_3),\,(w_2,w_4),\,(w_4,w_4)\}
\]

\begin{figure}[H]
\centering
\begin{tikzpicture}[node distance=2.8cm]
  \node[world] (w1) {$w_1$\\[-2pt]\tiny label};
  \node[world, right=of w1] (w3) {$w_3$\\[-2pt]\tiny artiste};
  \node[world, below=2cm of w1] (w2) {$w_2$\\[-2pt]\tiny fan};
  \node[world, below=2cm of w2] (w4) {$w_4$\\[-2pt]\tiny critique};

  \draw[arrow] (w1) -- (w3);
  \draw[arrow] (w1) -- (w2);
  \draw[arrow] (w2) -- (w4);
  \draw[arrow] (w3) edge[loop right] (w3);
  \draw[arrow] (w4) edge[loop right] (w4);
\end{tikzpicture}
\caption{Modèle de Kripke — domaine musical}
\label{fig:kripke-music}
\end{figure}

\noindent\textbf{Valuation $V$ :}
\begin{center}
\begin{tabular}{ll}
\toprule
Atome & Mondes vrais \\
\midrule
\code{album\_est\_platine}    & $\{w_1, w_3\}$ \\
\code{concert\_annule}        & $\{w_2, w_4\}$ \\
\code{artiste\_sous\_contrat} & $\{w_1\}$ \\
\code{chanson\_en\_top\_charts} & $\{w_1, w_2, w_3, w_4\}$ \\
\code{artiste\_independant}   & $\{w_3\}$ \\
\bottomrule
\end{tabular}
\end{center}

\subsection{Requêtes et résultats}

\begin{table}[H]
\centering
\caption{Évaluation des formules modales — domaine musical}
\begin{tabularx}{\linewidth}{clccX}
\toprule
\textbf{\#} & \textbf{Formule} & \textbf{Monde} & \textbf{Résultat} & \textbf{Justification} \\
\midrule
Q1 & $\nec\,\text{album\_platine}$ & $w_1$ & \textcolor{red}{\textbf{FAUX}}
   & $w_2 \in \mathrm{Succ}(w_1)$ mais album faux en $w_2$ \\
Q2 & $\pos\,\text{concert\_annule}$ & $w_1$ & \textcolor{codegreen}{\textbf{VRAI}}
   & concert vrai en $w_2 \in \mathrm{Succ}(w_1)$ \\
Q3 & $\nec\,\text{chanson\_top}$ & $w_2$ & \textcolor{codegreen}{\textbf{VRAI}}
   & $\mathrm{Succ}(w_2) = \{w_4\}$, chanson vrai en $w_4$ \\
Q4 & $\pos\,\text{artiste\_independant}$ & $w_1$ & \textcolor{codegreen}{\textbf{VRAI}}
   & indépendant vrai en $w_3 \in \mathrm{Succ}(w_1)$ \\
Q5 & $\nec(\text{album} \to \text{contrat})$ & $w_1$ & \textcolor{red}{\textbf{FAUX}}
   & en $w_3$ : album=V, contrat=F $\Rightarrow$ implication=F \\
\bottomrule
\end{tabularx}
\end{table}

\noindent\textbf{Propriétés de $R$ :} non réflexive ($w_2$ sans boucle),
non transitive, non symétrique, \textbf{sérielle} (tout monde a un successeur).

\section{Application — Système judiciaire}

\subsection{Modèle de Kripke}

\[
  W = \{w_1=\text{juge},\; w_2=\text{avocat},\; w_3=\text{procureur},\;
        w_4=\text{témoin},\; w_5=\text{accusé}\}
\]
\[
  R = \{(w_1,w_1),\,(w_1,w_2),\,(w_1,w_3),\,(w_2,w_4),\,(w_3,w_5),\,(w_5,w_5)\}
\]

\begin{figure}[H]
\centering
\begin{tikzpicture}[node distance=2.5cm]
  \node[world] (w1) {$w_1$\\[-2pt]\tiny juge};
  \node[world, right=of w1] (w2) {$w_2$\\[-2pt]\tiny avocat};
  \node[world, below=2.2cm of w1] (w3) {$w_3$\\[-2pt]\tiny procureur};
  \node[world, right=of w2] (w4) {$w_4$\\[-2pt]\tiny témoin};
  \node[world, right=of w3] (w5) {$w_5$\\[-2pt]\tiny accusé};

  \draw[arrow] (w1) -- (w2);
  \draw[arrow] (w1) -- (w3);
  \draw[arrow] (w2) -- (w4);
  \draw[arrow] (w3) -- (w5);
  \draw[arrow] (w1) edge[loop left] (w1);
  \draw[arrow] (w5) edge[loop right] (w5);
\end{tikzpicture}
\caption{Modèle de Kripke — domaine judiciaire}
\label{fig:kripke-justice}
\end{figure}

\noindent\textbf{Valuation $V$ :}
\begin{center}
\begin{tabular}{ll}
\toprule
Atome & Mondes vrais \\
\midrule
\code{accuse\_est\_coupable}        & $\{w_3, w_5\}$ \\
\code{accuse\_est\_innocent}        & $\{w_2, w_4\}$ \\
\code{preuve\_valide}               & $\{w_1, w_3\}$ \\
\code{temoin\_credible}             & $\{w_1, w_2, w_4\}$ \\
\code{circonstances\_attenuantes}   & $\{w_2, w_5\}$ \\
\bottomrule
\end{tabular}
\end{center}

\subsection{Requêtes et résultats}

\begin{table}[H]
\centering
\caption{Évaluation des formules modales — domaine judiciaire}
\begin{tabularx}{\linewidth}{clccX}
\toprule
\textbf{\#} & \textbf{Formule} & \textbf{Monde} & \textbf{Résultat} & \textbf{Justification} \\
\midrule
Q1 & $\nec\,\text{preuve\_valide}$ & $w_1$ & \textcolor{red}{\textbf{FAUX}}
   & preuve faux en $w_2 \in \mathrm{Succ}(w_1)$ \\
Q2 & $\pos\,\text{accuse\_innocent}$ & $w_1$ & \textcolor{codegreen}{\textbf{VRAI}}
   & innocent vrai en $w_2 \in \mathrm{Succ}(w_1)$ \\
Q3 & $\nec\,\text{accuse\_coupable}$ & $w_3$ & \textcolor{codegreen}{\textbf{VRAI}}
   & $\mathrm{Succ}(w_3) = \{w_5\}$, coupable vrai en $w_5$ \\
Q4 & $\pos\,\text{circ\_attenuantes}$ & $w_1$ & \textcolor{codegreen}{\textbf{VRAI}}
   & atténuantes vrai en $w_2 \in \mathrm{Succ}(w_1)$ \\
Q5 & $\nec\,\text{temoin\_credible}$ & $w_2$ & \textcolor{codegreen}{\textbf{VRAI}}
   & $\mathrm{Succ}(w_2) = \{w_4\}$, crédible vrai en $w_4$ \\
\bottomrule
\end{tabularx}
\end{table}

\noindent\textbf{Propriétés de $R$ :} réflexive en $w_1$ seulement,
non symétrique, non transitive, \textbf{non sérielle} ($w_4$ n'a aucun successeur).

\section{Implémentation avec TweetyProject}

L'évaluation modale est implémentée manuellement, TweetyProject étant utilisé
pour les structures de données (\code{FolAtom}, \code{Predicate},
\code{MlHerbrandInterpretation}).

\begin{lstlisting}[caption={Construction du modèle de Kripke en Java}]
// Atomes propositionnels
FolAtom albumPlatine = new FolAtom(new Predicate("album_est_platine"));

// Mondes (ensemble d'atomes vrais dans ce monde)
MlHerbrandInterpretation w1 = new MlHerbrandInterpretation(
    List.of(albumPlatine, contrat, chanson));
MlHerbrandInterpretation w2 = new MlHerbrandInterpretation(
    List.of(concertAnnule, chanson));

// Relation d'accessibilite
Set<Pair<Interpretation,Interpretation>> pairs = new HashSet<>();
pairs.add(new Pair<>(w1, w2));  // w1 -> w2
pairs.add(new Pair<>(w1, w3));  // w1 -> w3
pairs.add(new Pair<>(w3, w3));  // w3 boucle (reflexif local)
AccessibilityRelation R = new AccessibilityRelation(pairs);

// Evaluation de □album_platine depuis w1
boolean result = evaluate(w1, new Necessity(albumPlatine), R);
// → false : w2 est successeur de w1 et album_platine est faux en w2
\end{lstlisting}

\noindent L'évaluation récursive gère les cas :
\code{FolAtom} (consulte la valuation),
\code{Negation}, \code{Conjunction}, \code{Disjunction}, \code{Implication},
\code{Necessity} (itère sur tous les successeurs),
\code{Possibility} (cherche un successeur satisfaisant).
